\chapter{Funciones vectoriales}

\newpage
\section{Funciones vectoriales}

\subsection{Definición de funciones vectoriales}

\begin{definicion}[Función vectorial]
    Una \textbf{función vectorial} es una función cuya salida es un vector. En otras palabras, a cada valor de la variable independiente le corresponde un vector del espacio.
\end{definicion}

Se escribe como

\[
\mathbf{r}:D\subseteq\mathbb{R}\longrightarrow\mathbb{R}^{n},
\]

donde:

\begin{itemize}
    \item $D$ es el dominio de la función.
    \item $\mathbb{R}^{n}$ es el espacio donde se encuentran los vectores de salida.
\end{itemize}

Una función vectorial puede expresarse mediante sus funciones componentes:

\[
\mathbf{r}(t)=\left(x(t),\,y(t),\,z(t)\right),
\]

o, de forma equivalente,

\[
\mathbf{r}(t)=x(t)\mathbf{i}+y(t)\mathbf{j}+z(t)\mathbf{k}.
\]

En dos dimensiones basta considerar las dos primeras componentes.


\begin{ejercicio}[evaluar función vectorial]
Dada

\[
\mathbf{r}(t)=\left(t+1,2t-3\right),
\]

calcule $\mathbf{r}(2)$.

\textbf{Solución}

\[
\mathbf{r}(2)
=
(3,1).
\]
\end{ejercicio}

\newpage

\subsection{Interpretación geométrica}

Cada valor del parámetro $t$ determina un punto del espacio.

Cuando el parámetro varía continuamente, el extremo del vector recorre una trayectoria denominada \textbf{curva parametrizada}.

\[
t
\quad\longrightarrow\quad
(x(t),y(t),z(t))
\]

Es importante notar que la variable independiente suele llamarse \emph{parámetro}, ya que describe el movimiento sobre la curva y no necesariamente representa una coordenada espacial.

\begin{ejemplo}[Recta]
Sea

\[
\mathbf{r}(t)=\left(2+t,\;1+3t\right).
\]

Calculamos algunos puntos.

\[
\begin{aligned}
t=0&\Rightarrow (2,1),\\
t=1&\Rightarrow (3,4),\\
t=2&\Rightarrow (4,7).
\end{aligned}
\]

Estos puntos pertenecen a una recta del plano.
\end{ejemplo}


\begin{ejemplo}[Circunferencia]

Considere la función

\[
\mathbf{r}(t)=\left(\cos t,\sin t\right).
\]

Algunos valores son

\[
\begin{array}{c|c}
t & \mathbf{r}(t)\\
\hline
0&(1,0)\\
\dfrac{\pi}{2}&(0,1)\\
\pi&(-1,0)\\
\dfrac{3\pi}{2}&(0,-1)
\end{array}
\]

Los puntos obtenidos recorren la circunferencia unitaria.
\end{ejemplo}

\newpage

\begin{ejemplo}[Parábola]

Sea

\[
\mathbf{r}(t)=\left(t,t^{2}\right).
\]

Para algunos valores del parámetro:

\[
\begin{array}{c|c}
t & \mathbf{r}(t)\\
\hline
-2&(-2,4)\\
-1&(-1,1)\\
0&(0,0)\\
1&(1,1)\\
2&(2,4)
\end{array}
\]

La trayectoria corresponde a la parábola

\[
y=x^{2}.
\]
\end{ejemplo}

\begin{ejemplo}[Hélice]
Considere
\[
\mathbf{r}(t)=\left(\cos t,\sin t,t\right).
\]
En este caso:
\begin{itemize}
    \item las componentes $x$ e $y$ describen una circunferencia;
    \item la componente $z$ aumenta linealmente;
    \item la trayectoria es una hélice.
\end{itemize}
\end{ejemplo}

\newpage

\subsection{Operaciones con funciones vectoriales}

\begin{definicion}[Suma]
Sean:
\begin{itemize}
    \item $\mathbf{u}(t)=\left(u_{1}(t),u_{2}(t)\right)$
    \item $\mathbf{v}(t)=\left(v_{1}(t),v_{2}(t)\right)$
\end{itemize} 
Entonces:
\[
\mathbf{u}(t)+\mathbf{v}(t)
=
\left(
u_{1}+v_{1},
u_{2}+v_{2}
\right).
\]
\end{definicion}

\begin{definicion}[Multiplicación por un escalar]
Sean $\mathbf{u}(t)=\left(u_{1}(t),u_{2}(t)\right)$ y $k \in \mathbb{R}$. Entonces:
\[
k\mathbf{u}(t)
=
\left(
ku_{1},
ku_{2}
\right).
\]
\end{definicion}

% \newpage
\section{Dominio de funciones vectoriales}

\begin{definicion}[Dominio de una función vectorial]
\begin{itemize}
    \item El \textbf{dominio} de una función vectorial es el conjunto de todos los valores del parámetro para los cuales \textbf{todas} sus funciones componentes están definidas.
    \item Si:
        \[
        \mathbf{r}(t)=\left(f(t),g(t),h(t)\right),
        \]
        , entonces su dominio es:
        \[
        \operatorname{Dom}(\mathbf{r})
        =
        \operatorname{Dom}(f)
        \cap
        \operatorname{Dom}(g)
        \cap
        \operatorname{Dom}(h).
        \]
        \item Es decir, el dominio de la función vectorial corresponde a la \textbf{intersección de los dominios} de cada una de sus componentes.
\end{itemize}
\end{definicion}


\subsection{¿Por qué se utiliza la intersección?}

Para que la función vectorial esté definida en un determinado valor de $t$, es necesario que \textbf{todas} sus componentes puedan evaluarse simultáneamente.

Si una sola componente no está definida para un cierto valor del parámetro, entonces la función vectorial completa tampoco está definida en dicho valor.

\begin{ejemplo}[]

Considere

\[
\mathbf{r}(t)=
\left(
\sqrt{t},
\ln t,
t^2
\right).
\]


Analizando cada componente:

\begin{itemize}
    \item $\sqrt{t}$ exige
    \[
    t\ge0.
    \]

    \item $\ln t$ exige
    \[
    t>0.
    \]

    \item $t^2$ está definida para todo número real.
\end{itemize}

Por lo tanto,

\[
[0,\infty)
\cap
(0,\infty)
\cap
\mathbb{R}
=
(0,\infty).
\]

Así,

\[
\boxed{\operatorname{Dom}(\mathbf{r})=(0,\infty).}
\]

\end{ejemplo}

\textbf{Procedimiento para determinar el dominio}\\

Para encontrar el dominio de una función vectorial es recomendable seguir los siguientes pasos:

\begin{enumerate}
    \item Identificar cada función componente.
    \item Determinar el dominio de cada componente por separado.
    \item Calcular la intersección de todos los dominios.
\end{enumerate}

Este procedimiento garantiza que todas las componentes puedan evaluarse simultáneamente.

\newpage 

\begin{ejercicio}[]

Determine el dominio de

\[
\mathbf{r}(t)=
\left(
t^2,
e^t,
\sin t
\right).
\]

\textbf{Solución}

Cada componente está definida para todo número real.

\[
\operatorname{Dom}(t^2)=\mathbb{R},
\]

\[
\operatorname{Dom}(e^t)=\mathbb{R},
\]

\[
\operatorname{Dom}(\sin t)=\mathbb{R}.
\]

Por lo tanto,

\[
\boxed{\operatorname{Dom}(\mathbf{r})=\mathbb{R}.}
\]

\end{ejercicio}

\begin{ejercicio}[]

Determine el dominio de

\[
\mathbf{r}(t)=
\left(
\sqrt{t},
\frac1{t-1}
\right).
\]

\textbf{Solución}

Para la primera componente,

\[
t\ge0.
\]

Para la segunda,

\[
t\neq1.
\]

La intersección es

\[
[0,\infty)\setminus\{1\}.
\]

Por tanto,

\[
\boxed{\operatorname{Dom}(\mathbf{r})=[0,\infty)\setminus\{1\}.}
\]

\end{ejercicio}

\begin{ejercicio}[]

Determine el dominio de

\[
\mathbf{r}(t)=
\left(
\ln(t-2),
\sqrt{5-t}
\right).
\]

\textbf{Solución}

Primera componente:

\[
t>2.
\]

Segunda componente:

\[
5-t\ge0
\quad\Longrightarrow\quad
t\le5.
\]

La intersección es

\[
(2,5].
\]

Luego,

\[
\boxed{\operatorname{Dom}(\mathbf{r})=(2,5].}
\]

\end{ejercicio}

\begin{ejercicio}[]

Determine el dominio de

\[
\mathbf{r}(t)=
\left(
\frac1{t^2-9},
\sqrt{t+4},
\ln(t)
\right).
\]

\textbf{Solución}

Analizando cada componente:

\[
\frac1{t^2-9}
\quad\Rightarrow\quad
t\neq-3,\;t\neq3.
\]

\[
\sqrt{t+4}
\quad\Rightarrow\quad
t\ge-4.
\]

\[
\ln(t)
\quad\Rightarrow\quad
t>0.
\]

La condición más restrictiva es

\[
t>0,
\]

pero además debe cumplirse

\[
t\neq3.
\]

Así,

\[
\boxed{\operatorname{Dom}(\mathbf{r})=(0,\infty)\setminus\{3\}.}
\]
\end{ejercicio}

\section*{Ejercicios Propuestos}

Determine el dominio de las siguientes funciones vectoriales.

\begin{enumerate}
    \item
    \[
    \mathbf{r}(t)=
    \left(
    \sqrt{t},
    t^3,
    \cos t
    \right).
    \]

    \item
    \[
    \mathbf{r}(t)=
    \left(
    \ln(t+1),
    \sqrt{4-t}
    \right).
    \]

    \item
    \[
    \mathbf{r}(t)=
    \left(
    \frac1{t},
    \frac1{t-2}
    \right).
    \]

    \item
    \[
    \mathbf{r}(t)=
    \left(
    \sqrt{t-5},
    \ln(t-1),
    \frac1{t+2}
    \right).
    \]

    \item
    \[
    \mathbf{r}(t)=
    \left(
    e^t,
    \sqrt{9-t^2},
    \tan t
    \right).
    \]
\end{enumerate}


\newpage
\section{Límites de funciones vectoriales}

El concepto de límite para funciones vectoriales es una extensión natural del límite de funciones reales. La diferencia es que ahora la función no devuelve un número, sino un vector.

La idea fundamental es la siguiente:

\begin{quote}
\centering
\textit{
``Cuando el parámetro $t$ se aproxima a un valor $a$, \\
el vector $\mathbf{r}(t)$ se aproxima a un vector fijo $\mathbf{L}$.''
}
\end{quote}

En términos geométricos, el punto que describe la función vectorial se acerca a una posición determinada en el espacio.

El límite se calcula componente a componente. Si $\mathbf{r}(t)=\left(f(t),g(t),h(t)\right),$ entonces:

\[
\lim_{t\to a}\mathbf{r}(t)
=
\left(
\lim_{t\to a}f(t),
\lim_{t\to a}g(t),
\lim_{t\to a}h(t)
\right),
\]
siempre que existan los límites de cada componente.

\subsection{Definición}
\begin{definicion}[Límite de una función vectorial]
Sea:
\[\mathbf{r}(t)=\left(f(t),g(t),h(t)\right)\]
Decimos que:
\[\lim_{t\to a}\mathbf{r}(t)=\mathbf{L},\]
si existe el vector
\[\mathbf{L}=(L_1,L_2,L_3)\]
tal que:
\[
\lim_{t\to a}f(t)=L_1,\qquad
\lim_{t\to a}g(t)=L_2,\qquad
\lim_{t\to a}h(t)=L_3.
\]
En consecuencia,
\[
\boxed{
\lim_{t\to a}\mathbf{r}(t) =
\left(
\lim_{t\to a}f(t),
\lim_{t\to a}g(t),
\lim_{t\to a}h(t)
\right)
}
\]
siempre que existan todos los límites de las funciones componentes.
\end{definicion}

\newpage

\begin{ejercicio}[límite de una función vectorial]
Calcule:
\[ \lim_{t\to0} \left( \frac{\sin t}{t}, t^{2} \right) \]
\textbf{Solución}

Aplicando el límite componente a componente:

\begin{itemize}
    \item $ \lim_{t\to0} \dfrac{\sin t}{t}=1 $
    \item $ \lim_{t\to0}t^{2}=0 $
\end{itemize}

Por tanto:
\[ \lim_{t\to0} \left( \frac{\sin t}{t}, t^{2} \right)  = (1,0) \]
\end{ejercicio}

\subsection{Interpretación geométrica}

Una función vectorial puede interpretarse como el movimiento de una partícula. Si $ \mathbf{r}(t)= (x(t),y(t),z(t)) $, entonces cada valor de $t$ determina una posición en el espacio. Cuando $ t\rightarrow a $, la partícula se aproxima al punto $(L_1,L_2,L_3).$

Es importante notar que el límite puede existir incluso cuando la función no esté definida en $t=a$.

\section*{Propiedad Fundamental}

El límite de una función vectorial existe si y sólo si existen los límites de todas sus componentes.

Esta propiedad simplifica enormemente el cálculo de límites, pues basta trabajar con funciones reales, cuyos métodos ya son conocidos.

\newpage

\begin{ejemplo}[ ]
Calcule

\[
\lim_{t\to2}
(3t,\;t^2).
\]

\textbf{Solución}

Calculamos cada componente.

Primera componente:

\[
\lim_{t\to2}3t=6.
\]

Segunda componente:

\[
\lim_{t\to2}t^2=4.
\]

Por tanto,

\[
\boxed{
\lim_{t\to2}(3t,t^2)
=
(6,4).
}
\]
\end{ejemplo}

\begin{ejemplo}[ ]
Calcule
\[
\lim_{t\to0}
\left(
\frac{\sin t}{t},
e^t,
t^3
\right).
\]
\textbf{Solución}
Se calcula cada componente por separado.

\[
\lim_{t\to0}\frac{\sin t}{t}=1,
\]

\[
\lim_{t\to0}e^t=1,
\]

\[
\lim_{t\to0}t^3=0.
\]
Luego,
\[
\boxed{
(1,1,0).
}
\]
\end{ejemplo}

\newpage

\begin{ejemplo}[ ]

Calcule

\[
\lim_{t\to1}
\left(
\frac{t^2-1}{t-1},
\sqrt{t+3}
\right).
\]

\textbf{Solución}

Para la primera componente,

\[
\frac{t^2-1}{t-1}
=
t+1,
\qquad t\neq1.
\]

Entonces,

\[
\lim_{t\to1}\frac{t^2-1}{t-1}=2.
\]

Para la segunda,

\[
\lim_{t\to1}\sqrt{t+3}=2.
\]

Finalmente,

\[
\boxed{(2,2).}
\]
\end{ejemplo}

\begin{ejemplo}[ ]

Determine si existe

\[
\lim_{t\to0}
\left(
\frac1t,\;
\sin t
\right).
\]

\textbf{Solución}

La primera componente no posee límite cuando

\[
t\rightarrow0.
\]

Aunque

\[
\lim_{t\to0}\sin t=0,
\]

el límite de la función vectorial no existe.

\[
\boxed{
\lim_{t\to0}
\left(
\frac1t,\sin t
\right)
\text{ no existe.}
}
\]
\end{ejemplo}


\begin{teorema}[\\Propiedades de los Límites de funciones vectoriales]
Sean $ \mathbf{r}(t), \; \mathbf{s}(t) $ funciones vectoriales y $ c \in \mathbb{R} $ una constante. Si existen los límites correspondientes, entonces:

\begin{enumerate}
    \item \textbf{Suma}:
        \[
        \lim_{t\to a}
        \left(
        \mathbf{r}(t)+\mathbf{s}(t)
        \right)
        =
        \lim_{t\to a}\mathbf{r}(t)
        +
        \lim_{t\to a}\mathbf{s}(t).
        \]
    \item \textbf{Multiplicación por un escalar}:
        \[
        \lim_{t\to a}
        c\mathbf{r}(t)
        =
        c
        \lim_{t\to a}\mathbf{r}(t).
        \]
    \item \textbf{Producto punto}: Si $ \mathbf{r}(t) \cdot \mathbf{s}(t) $
        está definido,
        \[
        \lim_{t\to a}
        \left(
        \mathbf{r}\cdot\mathbf{s}
        \right)
        =
        \left(
        \lim\mathbf{r}
        \right)
        \cdot
        \left(
        \lim\mathbf{s}
        \right).
        \]
    \item \textbf{Producto cruz}: En $\mathbb{R}^3$,
        \[
        \lim_{t\to a}
        \left(
        \mathbf{r}\times\mathbf{s}
        \right)
        =
        \left(
        \lim\mathbf{r}
        \right)
        \times
        \left(
        \lim\mathbf{s}
        \right).
        \]
\end{enumerate}
\end{teorema}

\section*{Relación entre Límite y Continuidad}

Una función vectorial es continua en $t=a$ si

\[
\lim_{t\to a}\mathbf{r}(t)
=
\mathbf{r}(a).
\]

Como consecuencia, una función vectorial será continua cuando todas sus componentes sean funciones continuas.

\section*{Errores Frecuentes}

\begin{enumerate}
    \item Intentar calcular el límite del vector completo sin analizar cada componente.

    \item Pensar que basta con que una componente tenga límite para que exista el límite vectorial.

    \item Olvidar simplificar expresiones indeterminadas antes de sustituir.

    \item Confundir el punto al que tiende el parámetro con el vector límite.
\end{enumerate}

\section*{Ejercicios Propuestos}

Calcule los siguientes límites.

\begin{enumerate}

\item

\[
\lim_{t\to3}
(t^2,\;2t,\;5).
\]

\item

\[
\lim_{t\to0}
(\sin t,\cos t).
\]

\item

\[
\lim_{t\to1}
\left(
\frac{t^2-1}{t-1},
t^3
\right).
\]

\item

\[
\lim_{t\to0}
\left(
\frac{\sin t}{t},
\frac{1-\cos t}{t}
\right).
\]

\item

\[
\lim_{t\to4}
(\sqrt{t},\ln t).
\]

\item

\[
\lim_{t\to0}
\left(
\frac1t,
t^2
\right).
\]

Indique si existe o no el límite.

\item

Sean

\[
\mathbf{u}(t)=
(t,2t)
\]

y

\[
\mathbf{v}(t)=
(t^2,1).
\]

Calcule

\[
\lim_{t\to1}
(\mathbf{u}(t)+\mathbf{v}(t)).
\]

\item

Sean

\[
\mathbf{u}(t)=
(t,1,t^2),
\qquad
\mathbf{v}(t)=
(2,t,1).
\]

Calcule

\[
\lim_{t\to2}
(\mathbf{u}(t)\cdot\mathbf{v}(t)).
\]

\end{enumerate}

\section*{Resumen}

Sea

\[
\mathbf{r}(t)=
(f(t),g(t),h(t)).
\]

El límite se calcula componente a componente:

\[
\boxed{
\lim_{t\to a}\mathbf{r}(t)
=
\left(
\lim_{t\to a}f(t),
\lim_{t\to a}g(t),
\lim_{t\to a}h(t)
\right).
}
\]

El límite existe únicamente cuando existen los límites de todas las componentes.

\bigskip

\begin{center}
\begin{tabular}{|c|c|}
\hline
Operación & Regla\\
\hline
Suma & Se calcula componente a componente\\
\hline
Escalar por vector & El escalar multiplica cada componente\\
\hline
Producto punto & El límite puede entrar al producto\\
\hline
Producto cruz & El límite puede entrar al producto\\
\hline
Continuidad & $\displaystyle \lim_{t\to a}\mathbf{r}(t)=\mathbf{r}(a)$\\
\hline
\end{tabular}
\end{center}

% \newpage
\section{Derivadas de funciones vectoriales}

La derivada también se calcula componente a componente.

Si

\[
\mathbf{r}(t)=\left(x(t),y(t),z(t)\right),
\]

entonces

\[
\mathbf{r}'(t)=
\left(
x'(t),
y'(t),
z'(t)
\right).
\]

Geométricamente, la derivada representa el vector tangente a la curva.

\begin{ejemplo}[Derivada de una función vectorial]

Sea

\[
\mathbf{r}(t)=
\left(
t^{2},
\sin t,
e^{t}
\right).
\]

Entonces

\[
\mathbf{r}'(t)=
\left(
2t,
\cos t,
e^{t}
\right).
\]
\end{ejemplo}

\newpage

\begin{ejercicio}[derivar función vectorial]
Determine la derivada de

\[
\mathbf{r}(t)=
\left(
t^{3},
\cos t,
\sqrt{t}
\right).
\]

\textbf{Solución}

\[
\mathbf{r}'(t)=
\left(
3t^{2},
-\sin t,
\frac{1}{2\sqrt{t}}
\right).
\]
\end{ejercicio}

% \newpage
\section{Integrales de funciones vectoriales}

% \newpage
\section{Continuidad en funciones vectoriales}
Una función vectorial es continua si todas sus funciones componentes son continuas.


% \newpage
\section{Aplicaciones de funciones vectoriales}

% \newpage
\section{Resumen}

\begin{center}
\begin{tabular}{|c|c|}
\hline
Función vectorial & $\mathbf{r}:D\subseteq\mathbb{R}\rightarrow\mathbb{R}^{n}$\\
\hline
Forma general & $\mathbf{r}(t)=(x(t),y(t),z(t))$\\
\hline
Límite & Componente a componente\\
\hline
Continuidad & Todas las componentes continuas\\
\hline
Derivada & $\mathbf{r}'(t)=(x',y',z')$\\
\hline
Interpretación & Curva parametrizada\\
\hline
\end{tabular}
\end{center}

Para una función vectorial

\[
\mathbf{r}(t)=
\left(
f(t),g(t),h(t)
\right),
\]

su dominio se obtiene mediante

\[
\boxed{
\operatorname{Dom}(\mathbf{r})
=
\operatorname{Dom}(f)
\cap
\operatorname{Dom}(g)
\cap
\operatorname{Dom}(h).
}
\]

Es decir, el dominio está formado por aquellos valores del parámetro para los cuales \textbf{todas} las componentes están definidas simultáneamente.



\section{Ejercicios Propuestos}

\begin{enumerate}
\item Calcule $\mathbf{r}(3)$ si

\[
\mathbf{r}(t)=\left(2t-1,t^{2}+1\right).
\]

\item Determine la derivada de

\[
\mathbf{r}(t)=
\left(
e^{t},
\ln t,
t^{4}
\right).
\]

\item Calcule

\[
\lim_{t\to1}
\left(
t^{2},
\frac{t^{3}-1}{t-1}
\right).
\]

\item Verifique que

\[
\mathbf{r}(t)=
(\cos t,\sin t)
\]

describe una circunferencia de radio 1.

\item Construya una tabla de valores para

\[
\mathbf{r}(t)=
(t,t^{2})
\]

utilizando

\[
t=-2,-1,0,1,2.
\]

Grafique los puntos e identifique la curva obtenida.

\item Determine el vector tangente de

\[
\mathbf{r}(t)=
(t^{2},t^{3})
\]

en $t=1$.

\item Considere

\[
\mathbf{r}(t)=
(\cos t,\sin t,t).
\]

Describa verbalmente la trayectoria que genera esta función.

\item Escriba una función vectorial que describa una recta que pase por el punto $(1,2)$ y tenga dirección $(3,-1)$.
\end{enumerate}
